\section*{Exercise (iv)}

We assume $V^0(t,x) = x$ and then show $V^0(t,x)$ satisfies the PDE from Proposition~3.3.
\\
\\
According to Proposition~3.3, assuming $V^0(t,x)$ is sufficiently differentiable and $V^1 \equiv 0$, the market reserve must satisfy three conditions:
\begin{align}
  V_t^0(t, x) & = rV^0(t, x) - b^0(t, x) - \mu(t)R^{01}(t, x) - \mathcal{D}_x V^0(t, x) \\
  V^0(t-, x)  & = \Delta B^0(t,x) + V^0\bigl(t, x - \Delta A^0(t,x)\bigr)               \\
  V^0(n, x)   & = 0
\end{align}

We start by verifying the simplest conditions for $V^0(t,x) = x$.
\\
\\
For (2): Since the contract features no discrete jumps in benefits or account balance, $\Delta B^0 = \Delta A^0 = 0$. The condition trivially reduces to $x = x$, which is satisfied.
\\
\\
For (3): At the terminal age $n=111$, the basis reserve computed in Exercise~(i) satisfies $\tilde{V}^0(111)=0$. This forces the annuity payout rate $X(t)/\tilde{V}^0(t) \to \infty$, draining the account so that $X(111) = 0$. Thus, $V^0(111,0) = 0$, satisfying the terminal condition.
\\
\\
For (1): As a candidate for the reserve, let $V^0(t,x) = x$. We calculate its partial derivatives:
\[
  V_t^0(t, x) = 0, \qquad V_x^0(t, x) = 1, \qquad V_{xx}^0(t, x) = 0.
\]

$\hspace{0.5cm}$

Substituting these derivatives into $\mathcal{D}_x V^0(t, x)$ yields:
\begin{align*}
  \mathcal{D}_x V^0(t, x) & = V_x^0(t, x) \bigl(rx - a^0(t,x) - \mu(t)\varrho^{01}(t,x)\bigr) + \frac{1}{2}V_{xx}^0(t, x)\pi^2(t, x)\sigma^2 x^2 \\
                          & = 1 \cdot \bigl(rx - a^0(t,x) - \mu(t)\varrho^{01}(t,x)\bigr) + 0                                                    \\
                          & = rx - a^0(t,x) - \mu(t)\varrho^{01}(t,x).
\end{align*}
Furthermore, the sum at risk is $R^{01}(t,x) = b^{01}(t,x) - x$.

By the definition of a unit-linked contract, the account perfectly matches the benefits:
\[
  a^0(t,x) = b^0(t,x) \quad \text{and} \quad a^{01}(t,x) = b^{01}(t,x).
\]
Consequently, since the account terminates upon death, the internal account value jumps to $\chi^1(t,x) = 0$. The total account jump upon death therefore exactly matches the sum at risk:
\[
  \varrho^{01}(t,x) = a^{01}(t,x) + \chi^1(t,x) - x = b^{01}(t,x) + 0 - x = R^{01}(t,x).
\]

Inserting these equivalences into the right-hand side of (1), it simplifies analytically to zero:
\begin{align*}
  V_t^0(t,x) & = r x - b^0(t,x) - \mu(t) R^{01}(t, x) - \mathcal{D}_x V^0(t, x)                                     \\
             & = r x - a^0(t,x) - \mu(t) \varrho^{01}(t, x) - \bigl(rx - a^0(t,x) - \mu(t) \varrho^{01}(t, x)\bigr) \\
             & = 0.
\end{align*}

Since $V_t^0 = 0$, the PDE evaluates perfectly. Having satisfied all three conditions of Proposition~3.3, the verification theorem guarantees our candidate is the unique prospective market reserve. Thus, $V^0(t,x) = x$.

